\section{How to Read Core 1.3}

Core 1.3 is designed for several readers at once, and the structure of the volume reflects that
fact explicitly rather than forcing all readers through one rhetorical funnel.
The monograph is meant to be legible to editors and reviewers who need a clean scientific
surface, to formal readers who need exact definitions and theorem dependencies, to broad
scientific readers who need patient explanation before abstraction, and to domain readers who
need to understand what the framework does and does not yet prove about later projections.

\subsection{Reader tiers}

\textbf{Editors and scientific reviewers.}
The fastest high-trust route is to read the abstract, the present reader guide, the core
results chapters, the discussion and falsifiability chapters, and then the 1.3 source-audit
chapter near the end of the main body. That path exposes the scientific promise of the volume,
the structure of the theorem-bearing claims, and the reason why the 1.3 release is more honest
than a compressed packet that hides unresolved inheritance behind style.

\textbf{Formal readers.}
If the main concern is theorem discipline, the natural path is through the axiomatics,
operator semantics, theorem corpus, collapse and rebirth chapters, and the appendices.
This route keeps definitions, hypotheses, and derived consequences visible in the same volume,
so the reader does not need to shuttle between a small outward article and a long internal
registry simply to check whether a theorem uses what it says it uses.

\textbf{Broad scientific readers.}
Readers coming from systems theory, complexity science, or adjacent foundational work can read
the introduction, background, model, results, discussion, and conclusion first, then use the
reader guide sections of the later 1.3 material to understand which parts of the inherited OC
corpus are presented as stable core, which parts are bounded, and which parts remain
forward-looking rather than fully defended.

\textbf{Domain readers.}
Later chapters on levels, spaces, cycles, experiments, falsifiability, branching topology,
disciplines, modules, and the observed-universe contour remain part of the volume because the
framework is not only a slogan about collapse. It is a structural program whose scientific
meaning depends on how its formal core projects into actual model classes and empirical or
operational test lanes.

\subsection{Why the monograph remains necessary}

The short journal-core article extracted from Core 1.3 has a different purpose.
It is meant to present one bounded theorem program in a form that a journal editor or a hostile
reviewer can inspect quickly. That is valuable, but it is not sufficient as the canonical
scientific home of the framework. A foundational theory becomes less trustworthy, not more,
when the master document is replaced by a short packet that omits the didactic ladder, the
axiomatic shell, the inherited chapter architecture, and the boundary between what is structural
core and what is later projection.

For that reason the monograph remains the primary source. It gives the reader the whole
scientific object: the vocabulary, the theory, the theorems, the proofs, the illustrations, the
limitations, and the 1.3-specific revalidation apparatus that explains how the present release
relates to the inherited Core 1.2 corpus and to later OC/ESTRA work.

\subsection{What is new in the 1.3 release logic}

The 1.3 release is not just ``Core 1.2 plus a new PDF''.
It adds an explicit source-first publication discipline:

\begin{itemize}
\item the full LaTeX corpus remains the canonical authoring substrate rather than being treated as
      a raw archive behind a much smaller outward article;
\item theorem-bearing claims are separated by fate, so the reader can see what survives as stated,
      what is revised and bounded, and what is kept outside the positive body of narrower derived
      artifacts;
\item ORCID-first scholarly identity replaces the previous email-forward publication shell;
\item the monograph and the journal-core article are split cleanly, so the owner and the reader
      know which artifact is the master scientific volume and which one is the condensed
      journal-facing extraction.
\end{itemize}

\subsection{How the book is organized}

The first large block of the volume preserves the inherited OC corpus: introduction, background,
formal model, results, discussion, conclusion, figures, boundaries, thresholds, the full K-level
hierarchy, operators, collapse and rebirth, branching topology, domain extensions, falsifiability,
and the modular architecture. Those chapters are not background noise; they are the formal and
didactic body that made Core 1.2 monograph-scale in the first place.

The second large block is the 1.3-specific repair and publication layer.
It explains the source-native witness discipline, the theorem-fate partition, the chronology of
what changed between the inherited release and the present source-first publication line, and the
relationship between the full monograph, the extracted journal-core article, and the broader
outbound system in Logion.

\subsection{A note on truthfulness}

This volume is intentionally structured to reduce two common failure modes.
The first is inflated rhetoric: calling something complete simply because it is large.
The second is misleading austerity: calling something rigorous simply because it is small.
Core 1.3 tries to avoid both.
It remains large enough to teach the theory properly and to carry its formal structure in one
place, while also making its boundary conditions explicit enough that narrower derived artifacts
can be judged without hidden borrowings.

\clearpage
